\begin{table}[t]
\centering
\small
\begin{tabular}{rccccc}
\hline
prime & $p \bmod 8$ & branch & channel & left inverse & congruence \\
\hline
32633 & 1 & \textbf{reversed} & selfdual & yes & holds \\
32647 & 7 & plain & selfdual & yes & holds \\
32653 & 5 & plain & selfdual & yes & holds \\
32687 & 7 & plain & selfdual & yes & holds \\
32693 & 5 & plain & selfdual & yes & holds \\
32707 & 3 & \textbf{reversed} & selfdual & yes & holds \\
32713 & 1 & plain & selfdual & yes & holds \\
32717 & 5 & plain & selfdual & yes & holds \\
32719 & 7 & plain & selfdual & yes & holds \\
32749 & 5 & plain & selfdual & yes & holds \\
32771 & 3 & \textbf{reversed} & selfdual & yes & holds \\
\hline
\end{tabular}
\caption{The orientation-normalised frame at eleven primes covering all four usable residue classes modulo $8$. Every prime lands on the self-dual channel. The three needing the reversed square-root branch are $32633$, $32707$ and $32771$, in classes $1$, $3$ and $3$ respectively --- so the branch requirement is not a function of the residue class. Verified by \texttt{scripts/verify\_orientation\_canonical\_independent.py}, which does not call the production selector.}
\label{tab:orientation}
\end{table}
